\section{Overview}\label{m:sec:overview}

This chapter introduces the stochastic process classes and analytical methods used
throughout the thesis.  Its organising distinction is between discrete-time and
continuous-time Markov chains.  In discrete time, the Galton--Watson process shows
how independence converts reproduction into functional iteration.  In continuous
time, exponential holding times and the generator convert event rates into
Kolmogorov equations, with birth--death and birth--death--catastrophe processes as
the principal count-valued examples.  The intended framework also includes
time-inhomogeneous chains and coupled ODE--CTMC systems.  The available sources do
not yet specify the logistic-speciation or coupled-system models, however, so the
corresponding subsections mark those gaps explicitly rather than filling them with
unsupported mathematics.
\begin{figure}[tbp]
  \centering
  \begin{tikzpicture}[
    font=\small,
    node distance=6mm,
    process/.style={draw=blue!60!black, fill=blue!7, rounded corners=3pt,
      minimum width=6.2cm, minimum height=9mm, align=center},
    method/.style={draw=teal!65!black, fill=teal!7, rounded corners=3pt,
      minimum width=6.2cm, minimum height=11mm, align=center},
    outcome/.style={draw=orange!70!black, fill=orange!9, rounded corners=3pt,
      minimum width=6.2cm, minimum height=11mm, align=center},
    flow/.style={-{Stealth[length=2.2mm]}, semithick, draw=gray!70!black}
  ]
    \node[process] (dt) {Discrete time: Galton--Watson chain $(Z_n)$};
    \node[process, right=14mm of dt] (ct)
      {Continuous time: birth--death and BDC chains $(X_t)$};

    \node[method, below=of dt] (dtm)
      {PGF composition and first-step analysis\\[-1pt]
       $\phi_{n+1}=\phi_n\!\circ\!\phi$};
    \node[method, below=of ct] (ctm)
      {Generator, backward equations and PGF PDEs\\[-1pt]
       $\partial_tG=q(z)\,\partial_zG$};

    \node[outcome, below=of dtm] (dto)
      {Survival recursion and late-time scale\\[-1pt]
       $S_n\sim A(p)(2p)^n$};
    \node[outcome, below=of ctm] (cto)
      {Extinction, conditional means and characteristics\\[-1pt]
       $\mathbb E[X_t\mid X_t>0]\to1/A_{\mathrm c}(p)$};

    \draw[flow] (dt) -- (dtm);
    \draw[flow] (dtm) -- (dto);
    \draw[flow] (ct) -- (ctm);
    \draw[flow] (ctm) -- (cto);
    \draw[flow, dashed, <->] (dto.east) -- node[above, font=\footnotesize,
      fill=white, inner sep=2pt] {comparison} (cto.west);
  \end{tikzpicture}
  \caption{The two analytical lanes developed in this chapter.  Discrete
  functional iteration leads to the constant $A(p)$, while continuous-time
  generator methods give the explicit comparison constant $A_{\mathrm c}(p)$
  and the PDEs treated by characteristics.}
  \label{m:fig:chapter-map}
\end{figure}

The second half of the chapter is organised by method rather than by process.
For discrete chains, probability generating functions turn independent sums into
products and branching into composition, while first-step analysis isolates one
transition and returns a self-consistent equation.  Together these methods lead to
the binary Galton--Watson survival recursion and the discrete constant $A(p)$.
Only the limiting conditional mean requires that constant here; its product,
series, bounds and Koenigs-function theory remain the subject of \ChA.

For continuous chains, the generator supplies forward and backward equations,
while generating functions compress an infinite system of state-probability
equations into a first-order partial differential equation.  The linear
birth--death process serves as the main example because the same model makes the
mean, extinction probability, moment calculations and conditioning explicit.  In
particular, it gives the continuous comparison constant
\[
  \Ac(p)=\frac{1-2p}{1-p},
\]
which can be derived in closed form and compared directly with the discrete
$A(p)$ without importing the latter's deeper theory.

The method of characteristics comes last, where it converts a first-order
generating-function PDE into ordinary differential equations along characteristic
curves.  The main text works through one absorption--death model whose solution can
also be obtained independently, so each characteristic step can be checked.  The
broader absorption catalogue is compressed into
\cref{m:app:extended-absorption}, preserving the extension without obscuring the
methodological spine of the chapter.

Two further appendices preserve detail that would otherwise overload this
introductory route.  Appendix~\ref{m:app:critical-gw} proves the critical
Galton--Watson survival and total-progeny asymptotics, while the following parts
record the limiting conditional variance, small-founder effects and the discrete
rupture calculation.  They support the claims used in the main text without
transferring the deeper theory of $A(p)$ from \ChA.

Absorption provides the common thread.  Extinction, killing and catastrophe all
remove mass from the transient state space, while conditioning asks what remains
visible before that absorption occurs.  This viewpoint connects elementary
first-step arguments to the spectral language of quasi-stationarity and to the
generating-function calculations used later in the thesis.  Standard background
for the continuous-time material is given by Norris \cite{Norris1997}, Karlin and
Taylor \cite{KarlinTaylor1975}, and Allen \cite{Allen2010}; the branching-process
background follows Harris \cite{Harris1963} and Athreya and Ney
\cite{Athreya1972}.
